%! AP Statistics — Unit 1, Lesson 1.1 — Slides (Beamer)
\documentclass[aspectratio=169, 11pt]{beamer}
\usepackage{apstats-beamer}

%=============================================================================
\begin{document}

% ─────────────────────────────────────────────────────────────────────────────
% SLIDE 2 — LESSON TITLE
% ─────────────────────────────────────────────────────────────────────────────
{
\setbeamercolor{background canvas}{bg=navy}
\begin{frame}[plain]
  \begin{tikzpicture}[remember picture, overlay]
    \fill[navylight] (current page.north west)
      rectangle ([yshift=-1.35cm]current page.north east);
  \end{tikzpicture}
  \vspace{2.8cm}
  \hspace{0.55cm}%
  \begin{minipage}{0.9\textwidth}
    {\color{goldacc}\bfseries\small LESSON 1.1}\\[0.5cm]
    {\color{white}\bfseries\LARGE Introducing Statistics:\\[0.25cm]
      What Can We Learn from Data?}\\[1.4cm]
    {\color{skymid}\small AP Statistics~$\cdot$~Unit 1: Exploring One-Variable Data}
  \end{minipage}
\end{frame}
}

% ─────────────────────────────────────────────────────────────────────────────
% SLIDE 3 — PRIMARY OBJECTIVE
% ─────────────────────────────────────────────────────────────────────────────
{
\setbeamercolor{background canvas}{bg=sky}
\begin{frame}[t, plain]
  \begin{tikzpicture}[remember picture, overlay]
    \fill[navy] (current page.north west)
      rectangle ([yshift=-0.25cm]current page.north east);
  \end{tikzpicture}
  \vspace{0.45cm}\par
  \hspace{0.55cm}%
  \begin{minipage}{0.9\textwidth}

    \sectionlabel{Primary Objective}

    \normalsize
    Students will understand that \textbf{statistics is the science of
    collecting, analyzing, and drawing conclusions from data.}
    Students will distinguish between populations and samples,
    identify individuals and variables in a dataset, and recognize
    why variation in data is central to statistical thinking
    (Big Idea: VAR\,---\,Skill 1).

    \vspace{0.6cm}

    \normalsize\textbf{\color{navy}By the end of this lesson, I will be able to:}

    \smallskip
    \begin{itemize}
      \setlength{\itemsep}{0.3em}\setlength{\topsep}{0.2em}
      \item Identify individuals, populations, samples, and variables in a
        statistical context.
      \item Explain why statistics is needed --- because individuals vary,
        and data captures that variation.
      \item Describe the four-step statistical problem-solving process.
    \end{itemize}

  \end{minipage}
\end{frame}
}

% ─────────────────────────────────────────────────────────────────────────────
% SLIDE 4 — HOOK
% ─────────────────────────────────────────────────────────────────────────────
{
\setbeamercolor{background canvas}{bg=hookbg}
\setbeamercolor{itemize item}{fg=white}
\begin{frame}[t, plain]
  \vspace{0.45cm}\par
  \hspace{0.55cm}%
  \begin{minipage}{0.9\textwidth}

    {\color{white}\bfseries\footnotesize\MakeUppercase{Hook}}\par\medskip

    {\color{white}\large\bfseries
      ``Students who sleep more score higher on tests.''}\\[0.25cm]
    {\color{white}\small\itshape
      Headline from a national education study ($n = 12{,}000$ students)}

    \vspace{0.5cm}
    {\color{white}\bfseries Discussion questions:}\par\smallskip
    \begin{itemize}
      \setlength{\itemsep}{0.4em}\setlength{\topsep}{0.1em}
      \item {\color{white}How do researchers know this?
        Did they study every student in the world?}
      \item {\color{white}Could this finding be wrong?}
      \item {\color{white}What information would you want before trusting
        this headline?}
    \end{itemize}

  \end{minipage}
\end{frame}
}

% ─────────────────────────────────────────────────────────────────────────────
% SLIDE 5 — VOCABULARY & KEY CONCEPTS
% ─────────────────────────────────────────────────────────────────────────────
{
\setbeamercolor{background canvas}{bg=greenbg}
\begin{frame}[t, plain]
  \begin{tikzpicture}[remember picture, overlay]
    \fill[greenacc] (current page.north west)
      rectangle ([yshift=-1.35cm]current page.north east);
    \node[anchor=west, text=white, font=\bfseries\large,
          inner xsep=0.55cm]
      at ([yshift=-0.68cm]current page.north west)
      {Vocabulary \& Key Concepts};
  \end{tikzpicture}
  \vspace{1.05cm}\par
  \hspace{0.4cm}%
  \begin{minipage}{0.93\textwidth}
    \vspace{1em}
    \footnotesize
    \begin{multicols}{2}
      \textbf{\color{greenacc}Statistics} --- the science of collecting,
        analyzing, and drawing conclusions from data\\[4pt]
      \textbf{\color{greenacc}Individual} --- a single person, animal, object,
        or event in a dataset\\[4pt]
      \textbf{\color{greenacc}Population} --- the entire group of individuals
        of interest\\[4pt]
      \textbf{\color{greenacc}Sample} --- a subset of the population from which
        data are actually collected

      \columnbreak

      \textbf{\color{greenacc}Census} --- data collected from every individual
        in the population\\[4pt]
      \textbf{\color{greenacc}Variability / Variation} --- natural differences
        in values across individuals; the reason statistics exists\\[4pt]
      \textbf{\color{greenacc}Data} --- recorded values for individuals;
        can be categorical or quantitative\\[4pt]
      \textbf{\color{greenacc}Anecdotal evidence} --- data from a few
        unrepresentative cases; not reliable for inference
    \end{multicols}
  \end{minipage}
\end{frame}
}

% ─────────────────────────────────────────────────────────────────────────────
% SLIDE 6 — THE STATISTICAL PROBLEM-SOLVING PROCESS
% ─────────────────────────────────────────────────────────────────────────────
{
\setbeamercolor{background canvas}{bg=white}
\begin{frame}[t, plain]
  \navyheader{The Statistical Problem-Solving Process}
  \hspace{0.55cm}%
  \begin{minipage}{0.9\textwidth}
    \vspace{1em}
    \small
    \begin{itemize}
      \setlength{\itemsep}{0.45em}\setlength{\topsep}{0em}
      \item \textbf{1.\; Ask a question}\\
        Pose a clear question about a population of interest
      \item \textbf{2.\; Collect data}\\
        Select a sample and measure variables on each individual
      \item \textbf{3.\; Analyze the data}\\
        Summarize and display; compute statistics
      \item \textbf{4.\; Interpret results}\\
        Draw conclusions; acknowledge what the data cannot tell you
    \end{itemize}
  \end{minipage}
\end{frame}
}

% ─────────────────────────────────────────────────────────────────────────────
% SLIDE 7 — KEY CONCEPTS: POPULATION VS. SAMPLE
% ─────────────────────────────────────────────────────────────────────────────
{
\setbeamercolor{background canvas}{bg=white}
\begin{frame}[t, plain]
  \navyheader{Key Concepts: Population vs.\ Sample}
  \hspace{0.4cm}%
  \begin{minipage}{0.93\textwidth}
    \vspace{0.3cm}
    \begin{columns}[T]

      \vspace{1em}
      \begin{column}{0.48\textwidth}
        \small
        \textbf{\color{navy}Example: Identifying the Components}\\[4pt]
        {\color{slate}\itshape A researcher surveys 200 randomly chosen
          adults about their exercise habits.}\\[6pt]
        \begin{itemize}
          \setlength{\itemsep}{0.3em}\setlength{\topsep}{0em}
          \item \textbf{Individual:}\; one adult
          \item \textbf{Population:}\; all adults
          \item \textbf{Sample:}\; the 200 randomly chosen adults
          \item \textbf{Variable:}\; exercise habits (e.g., days per week)
        \end{itemize}
        \vspace{6pt}
        \textbf{\color{navy}Why use a sample instead of a census?}\\[3pt]
        A census is too costly, time-consuming, or impossible for large
        populations. A well-chosen sample gives reliable estimates.
      \end{column}

      \begin{column}{0.48\textwidth}
        \small
        \textbf{\color{navy}Why Statistics Needs Variation}\\[4pt]
        {\itshape ``If every individual were identical, we would
        \textbf{not} need statistics --- there would be nothing different
        to explain or measure.''}\\[6pt]
        Variation is natural and expected. Statistical methods describe,
        summarize, and make sense of it.\\[8pt]
        \textbf{\color{navy}Don't confuse these:}
        \begin{itemize}
          \setlength{\itemsep}{0.3em}\setlength{\topsep}{3pt}
          \item \textbf{Sample} = the group of individuals selected
          \item \textbf{Data} = the values recorded from that group
        \end{itemize}
      \end{column}

    \end{columns}
  \end{minipage}
\end{frame}
}

% ─────────────────────────────────────────────────────────────────────────────
% SLIDE 8 — GROUP ACTIVITY
% ─────────────────────────────────────────────────────────────────────────────
{
\setbeamercolor{background canvas}{bg=redbg}
\setbeamercolor{enumerate item}{fg=redacc}
\begin{frame}[t, plain]
  \begin{tikzpicture}[remember picture, overlay]
    \fill[redacc] (current page.north west)
      rectangle ([yshift=-1.35cm]current page.north east);
    \node[anchor=west, text=white, font=\bfseries\large,
          inner xsep=0.55cm, inner ysep=0pt]
      at ([yshift=-0.675cm]current page.north west)
      {Group Activity};
  \end{tikzpicture}
  \vspace{1.05cm}\par
  \hspace{0.55cm}%
  \begin{minipage}{0.9\textwidth}
    \vspace{1em}
    \small
    \textbf{Scenario Analysis --- Identifying Statistical Components}\\[6pt]
    \begin{enumerate}
      \setlength{\itemsep}{0.4em}\setlength{\topsep}{0em}
      \item Read each of the 3 provided scenarios.
      \item For each: identify the individual, population, sample (if any),
        and one variable measured.
      \item Explain why a census is or is not feasible for each scenario.
      \item Share your reasoning with another pair.
    \end{enumerate}
    \vspace{0.6cm}
    {\color{redacc}\bfseries\small TIME: 8--10 minutes}
  \end{minipage}
\end{frame}
}

% ─────────────────────────────────────────────────────────────────────────────
% SLIDE 9 — EXIT TICKET
% ─────────────────────────────────────────────────────────────────────────────
{
\setbeamercolor{background canvas}{bg=navy}
\setbeamercolor{enumerate item}{fg=white}
\begin{frame}[t, plain]
  \begin{tikzpicture}[remember picture, overlay]
    \fill[navylight] (current page.north west)
      rectangle ([yshift=-1.35cm]current page.north east);
  \end{tikzpicture}
  \vspace{0.45cm}\par
  \hspace{0.55cm}%
  \begin{minipage}{0.9\textwidth}
    {\color{goldacc}\bfseries\footnotesize\MakeUppercase{Exit Ticket}}
    \vspace{0.55cm}\par
    \small
    \begin{enumerate}
      \setlength{\itemsep}{0.55em}\setlength{\topsep}{0em}
      \item {\color{white}A hospital tracks recovery times for 50 patients
        after knee surgery.\\
        $\rightarrow$ Identify the individual, population, and sample.}
      \item {\color{white}Name one variable in this study. Does it vary
        across individuals?}
      \item {\color{white}Why might researchers use a sample instead of all
        knee-surgery patients?}
    \end{enumerate}
  \end{minipage}
\end{frame}
}

% ─────────────────────────────────────────────────────────────────────────────
% SLIDE 10 — CLOSING & HOMEWORK
% ─────────────────────────────────────────────────────────────────────────────
{
\setbeamercolor{background canvas}{bg=navy}
\setbeamercolor{itemize item}{fg=white}
\begin{frame}[t, plain]
  \begin{tikzpicture}[remember picture, overlay]
    \fill[navylight] (current page.north west)
      rectangle ([yshift=-1.35cm]current page.north east);
  \end{tikzpicture}
  \vspace{0.45cm}\par
  \hspace{0.55cm}%
  \begin{minipage}{0.9\textwidth}
    {\color{goldacc}\bfseries\footnotesize\MakeUppercase{Closing \& Homework}}
    \vspace{0.45cm}\par
    \begin{columns}[T]

      \begin{column}{0.54\textwidth}
        \small
        {\color{goldacc}\bfseries Key Reminders}\\[5pt]
        \begin{itemize}
          \setlength{\itemsep}{0.4em}\setlength{\topsep}{0em}
          \item {\color{white}\textbf{Association $\neq$ causation.}
            {\color{skymid} A pattern in data does not mean one variable
            causes another.}}
          \item {\color{white}\textbf{Sample $\neq$ data.}
            {\color{skymid} A sample is the group of individuals; data are
            the values recorded from them.}}
          \item {\color{white}\textbf{Variation is informative, not a flaw.}
            {\color{skymid} It is expected and is exactly why we need
            statistical tools.}}
          \item {\color{white}\textbf{The 4-step process will spiral all year.}
            {\color{skymid} Internalize it now --- every unit builds on it.}}
        \end{itemize}
      \end{column}

      \begin{column}{0.42\textwidth}
        \small
        {\color{goldacc}\bfseries Homework}\\[5pt]
        {\color{skymid}Complete the independent practice problems on
        today's topic:}\\[5pt]
        \begin{itemize}
          \setlength{\itemsep}{0.35em}\setlength{\topsep}{0em}
          \item {\color{white}Identifying individuals, populations, and
            samples}
          \item {\color{white}Distinguishing samples from census situations}
          \item {\color{white}Identifying and describing variables in context}
        \end{itemize}
        \vspace{0.4cm}
        {\color{skymid}\itshape\small Due at the start of our next class.}
      \end{column}

    \end{columns}
  \end{minipage}
\end{frame}
}

\end{document}
